%% generated by make_refs.py from the citations in paper.aux; do not edit
\begin{thebibliography}{99}
\small
\bibitem{Moore64} Moore, E.~F. (1964). The firing squad synchronization problem. In: E.~F. Moore (ed.), \emph{Sequential Machines: Selected Papers}, pp.~213--214. Reading, MA: Addison-Wesley.
\bibitem{Goto62} Goto, E. (1962). A minimal time solution of the firing squad problem. Dittoed course notes for Applied Mathematics~298, Harvard University, pp.~52--59.
\bibitem{Waksman66} Waksman, A. (1966). An optimum solution to the firing squad synchronization problem. \emph{Inform.\ and Control}, 9, 66--78.
\bibitem{Mazoyer86} Mazoyer, J. (1988). An overview of the firing squad synchronization problem. In: C.~Choffrut (ed.), \emph{Automata Networks}, Lecture Notes in Comput.\ Sci. 316, pp.~82--94. Springer.
\bibitem{Umeo09} Umeo, H. (2009). Firing squad synchronization problem in cellular automata. In: R.~A. Meyers (ed.), \emph{Encyclopedia of Complexity and Systems Science}, pp.~3537--3574. Springer.
\bibitem{Correa17} Correa, T., Gustavo, B., Lemos, L., Settle, A. (2017). An overview of recent solutions to and lower bounds for the firing synchronization problem. arXiv:1701.01045.
\bibitem{Balzer67} Balzer, R. (1967). An 8-state minimal time solution to the firing squad synchronization problem. \emph{Inform.\ and Control}, 10, 22--42.
\bibitem{Gerken87} Gerken, H.-D. (1987). \emph{{\"U}ber Synchronisationsprobleme bei Zellularautomaten}. Diplomarbeit, Institut f{\"u}r Theoretische Informatik, Technische Universit{\"a}t Braunschweig.
\bibitem{Mazoyer87} Mazoyer, J. (1987). A six-state minimal time solution to the firing squad synchronization problem. \emph{Theoret.\ Comput.\ Sci.}, 50, 183--238.
\bibitem{Yunes93} Yun{\`e}s, J.-B. (1993). \emph{Synchronisation et automates cellulaires: la ligne de fusiliers}. PhD thesis, Universit{\'e} Paris~7. Available from HAL, tel-00139049.
\bibitem{Sanders94} Sanders, P. (1994). Massively parallel search for transition-tables of polyautomata. In: \emph{Parcella '94: Proceedings of the VI International Workshop on Parallel Processing by Cellular Automata and Arrays}, pp.~99--108. Berlin: Akademie Verlag.
\bibitem{Clergue18} Clergue, M., Verel, S., Formenti, E. (2018). An iterated local search to find many solutions of the 6-states firing squad synchronization problem. \emph{Appl.\ Soft Comput.}, 66, 449--461.
\bibitem{NguyenMaignan20} Nguyen, T.~T., Maignan, L. (2020). Exploring millions of 6-state FSSP solutions: the formal notion of local CA simulation. In: \emph{Cellular Automata and Discrete Complex Systems (AUTOMATA 2020)}, Lecture Notes in Comput.\ Sci. 12286, pp.~1--13. Springer.
\bibitem{UmeoYanagihara07} Umeo, H., Yanagihara, T. (2007). A smallest five-state solution to the firing squad synchronization problem. In: \emph{Machines, Computations, and Universality (MCU 2007)}, Lecture Notes in Comput.\ Sci. 4664, pp.~291--302. Springer.
\bibitem{Yunes08} Yun{\`e}s, J.-B. (2008). A 4-states algebraic solution to linear cellular automata synchronization. \emph{Inform.\ Process.\ Lett.}, 107, 71--75.
\bibitem{UmeoKamikawaYunes09} Umeo, H., Kamikawa, N., Yun{\`e}s, J.-B. (2009). A family of smallest symmetrical four-state firing squad synchronization protocols for ring arrays. \emph{Parallel Process.\ Lett.}, 19, 299--313.
\bibitem{Mazoyer96} Mazoyer, J. (1996). On optimal solutions to the firing squad synchronization problem. \emph{Theoret.\ Comput.\ Sci.}, 168, 367--404.
\bibitem{YamashitaEtAl14} Yamashita, K., Nishitani, Y., Hirose, S., Okawa, S., Osato, N. (2014). The firing squad synchronization problem with sub-generals. \emph{Inform.\ Process.\ Lett.}, 114, 60--65.
\bibitem{Kobayashi19} Kobayashi, K. (2019). Nonexistence of minimal-time solutions for some variations of the firing squad synchronization problem having simple geometric configurations. arXiv:1909.10125.
\bibitem{MazoyerTerrier99} Mazoyer, J., Terrier, V. (1999). Signals in one-dimensional cellular automata. \emph{Theoret.\ Comput.\ Sci.}, 217, 53--80.
\bibitem{BerthiaumeEtAl04} Berthiaume, A., Bittner, T., Perkovi{\'c}, L., Settle, A., Simon, J. (2004). Bounding the firing synchronization problem on a ring. \emph{Theoret.\ Comput.\ Sci.}, 320, 213--228.
\bibitem{DAntonioDelzanno04} D'Antonio, M., Delzanno, G. (2004). SAT-based analysis of cellular automata. In: \emph{Cellular Automata (ACRI 2004)}, Lecture Notes in Comput.\ Sci. 3305, pp.~745--754. Springer.
\bibitem{WetzlerHeuleHunt14} Wetzler, N., Heule, M.~J.~H., Hunt, W.~A., Jr. (2014). DRAT-trim: efficient checking and trimming using expressive clausal proofs. In: \emph{Theory and Applications of Satisfiability Testing (SAT 2014)}, Lecture Notes in Comput.\ Sci. 8561, pp.~422--429. Springer.
\bibitem{CruzFilipeEtAl17} Cruz-Filipe, L., Heule, M.~J.~H., Hunt, W.~A., Jr., Kaufmann, M., Schneider-Kamp, P. (2017). Efficient certified RAT verification. In: \emph{Automated Deduction -- CADE 26}, Lecture Notes in Comput.\ Sci. 10395, pp.~220--236. Springer.
\bibitem{HeuleKullmannMarek16} Heule, M.~J.~H., Kullmann, O., Marek, V.~W. (2016). Solving and verifying the Boolean Pythagorean triples problem via cube-and-conquer. In: \emph{Theory and Applications of Satisfiability Testing (SAT 2016)}, Lecture Notes in Comput.\ Sci. 9710, pp.~228--245. Springer.
\bibitem{Heule18} Heule, M.~J.~H. (2018). Schur number five. In: \emph{Proceedings of the Thirty-Second AAAI Conference on Artificial Intelligence (AAAI-18)}, pp.~6598--6606.
\bibitem{CaDiCaL2} Biere, A., Faller, T., Fazekas, K., Fleury, M., Froleyks, N., Pollitt, F. (2024). CaDiCaL~2.0. In: \emph{Computer Aided Verification (CAV 2024)}, Lecture Notes in Comput.\ Sci. 14681, pp.~133--152. Springer.
\bibitem{Kissat20} Biere, A., Fazekas, K., Fleury, M., Heisinger, M. (2020). CaDiCaL, Kissat, Paracooba, Plingeling and Treengeling entering the SAT Competition 2020. In: \emph{Proceedings of SAT Competition 2020: Solver and Benchmark Descriptions}, Department of Computer Science Report Series~B, B-2020-1, pp.~51--53. University of Helsinki.
\bibitem{Duprat97} Duprat, J. (1997). Firing Squad Synchronization Problem. Coq user contribution FiringSquad (formal proof of the correctness of Mazoyer's six-state solution), LIP, ENS Lyon. Archived at \url{https://github.com/rocq-archive/firing-squad}.
\bibitem{Wolfram83} Wolfram, S. (1983). Statistical mechanics of cellular automata. \emph{Rev.\ Modern Phys.}, 55, 601--644.
\end{thebibliography}
